
%% =====================================================================
%%  สไลด์ประกอบการสอน บทที่ 2 : ทฤษฎีเซต --- สัปดาห์ที่ 6
%%  หัวข้อ: เพาเวอร์เซต, ผลคูณคาร์ทีเซียน, จำนวนสมาชิก (Cardinality) และการประยุกต์
%%  อ้างอิงตรงตามเนื้อหาเอกสารประกอบการสอน (chapter02.tex)
%% =====================================================================

%% ====================== ปก ======================
\begin{frame}[plain]
  \maketitle
\end{frame}

%% ====================== ทบทวน ======================
\begin{frame}{ทบทวนสัปดาห์ที่ 5 : การดำเนินการและกฎพีชคณิตของเซต}
  \begin{itemize}
    \item \textbf{4 การดำเนินการพื้นฐาน}:
          $A \cup B$ (หรือ), \ $A \cap B$ (และ), \ $A - B$ (ผลต่าง $= A \cap B^c$), \ $A^c$ (คอมพลีเมนต์)
    \item \textbf{กฎพีชคณิตของเซต}: สมบัติการสลับที่ การเปลี่ยนหมู่ การแจกแจง เอกลักษณ์ และชดเชย
    \item \textbf{กฎเดอมอร์แกน (De Morgan's Laws)}:
          \[ (A \cup B)^c = A^c \cap B^c \qquad \text{และ} \qquad (A \cap B)^c = A^c \cup B^c \]
  \end{itemize}
  \vspace{0.2em}
  \begin{exampleblock}{สัปดาห์นี้ = สัปดาห์ปิดบทที่ 2 (แก่นแท้ของการสร้างและนับเซต)}
    \begin{itemize}\small
      \item \textbf{เพาเวอร์เซต} $\mathcal{P}(A)$ : รวบรวมเซตย่อยทั้งหมด และการเติบโต $2^{|A|}$
      \item \textbf{ผลคูณคาร์ทีเซียน} $A \times B$ : การจับคู่สร้างคู่อันดับ รากฐานของฐานข้อมูลและฟังก์ชัน
      \item \textbf{จำนวนสมาชิก (Cardinality)} : หลักการรวม-การแยก (Inclusion-Exclusion) ทั้ง 2 และ 3 เซต
      \item \textbf{เซตในคอมพิวเตอร์} : Bit Vector, SQL CROSS JOIN, Hash Set $O(1)$ และเซตอนันต์
    \end{itemize}
  \end{exampleblock}
\end{frame}

%% ====================== วัตถุประสงค์ ======================
\begin{frame}{วัตถุประสงค์การเรียนรู้ (สัปดาห์ที่ 6)}
  เมื่อเรียนจบสัปดาห์นี้ นักศึกษาจะสามารถ
  \begin{enumerate}
    \item หา \textbf{เพาเวอร์เซต} $\mathcal{P}(A)$ อธิบายที่มาของสูตร $|\mathcal{P}(A)| = 2^{|A|}$ และเชื่อมโยงกับการเติบโตเชิงเลขชี้กำลังและบิตมาสก์ได้
    \item หา \textbf{ผลคูณคาร์ทีเซียน} $A \times B$ เขียนตารางคาร์ทีเซียน ใช้สมบัติการกระจาย และคำนวณจำนวนสมาชิก $|A \times B| = |A| \cdot |B|$ ได้
    \item คำนวณ \textbf{จำนวนสมาชิกของเซต (Cardinality)} โดยใช้ \textbf{หลักการรวม-การแยก (Inclusion-Exclusion Principle)} สำหรับ 2 และ 3 เซต พร้อมตรวจความสอดคล้องของข้อมูลได้
    \item อธิบายการแทนเซตด้วย \textbf{Bit Vector} การทำงานของ \textbf{SQL / Hash Set} และเข้าใจแนวคิด \textbf{เซตนับได้/นับไม่ได้} ได้
  \end{enumerate}
\end{frame}

%% ====================== หัวข้อ ======================
\begin{frame}{หัวข้อการบรรยายประจำสัปดาห์}
  \tableofcontents
\end{frame}


%% =====================================================================
\section{เพาเวอร์เซต (Power Set)}
%% =====================================================================

\begin{frame}{เพาเวอร์เซต : นิยามและความหมาย}
  \begin{block}{นิยาม : เพาเวอร์เซต (Power Set)}
    \textbf{เพาเวอร์เซต} ของเซต $A$ เขียนแทนด้วยสัญลักษณ์ $\mathcal{P}(A)$ หรือ $2^A$ คือ
    \textbf{เซตที่ประกอบด้วยเซตย่อยทั้งหมด} ของ $A$ นั่นคือ
    \[ \mathcal{P}(A) = \{B \mid B \subseteq A\} \]
  \end{block}
  \vspace{0.2em}
  \textbf{ลักษณะสำคัญของเพาเวอร์เซต}:
  \begin{itemize}
    \item เพาเวอร์เซตเป็น \textbf{``เซตของเซต''} เพราะสมาชิกทุกตัวของมันคือ ``เซตย่อย'' ของ $A$
    \item สำหรับทุกเซต $A$ จะได้ว่า $\emptyset \in \mathcal{P}(A)$ และ $A \in \mathcal{P}(A)$ เสมอ
    \item $\emptyset \subseteq \mathcal{P}(A)$ เสมอ (เพราะเซตว่างเป็นเซตย่อยของทุกเซต)
  \end{itemize}
  \vspace{0.2em}
  \begin{alertblock}{ข้อควรระวังสำคัญมาก!}
    $\mathcal{P}(\emptyset) = \{\emptyset\}$ \quad \textbf{ไม่ใช่เซตว่าง!} \\
    $\mathcal{P}(\emptyset)$ มีสมาชิก 1 ตัว คือ เซตว่าง ($\emptyset$) ดังนั้น $|\mathcal{P}(\emptyset)| = 2^0 = 1$
  \end{alertblock}
\end{frame}

\begin{frame}{ตัวอย่างการแจกแจงเพาเวอร์เซต}
  ลองไล่เขียนเพาเวอร์เซตจากขนาดเล็กไปใหญ่ โดยเริ่มจากเซตย่อยขนาด 0, 1, 2, \ldots
  \begin{center}\small
  \begin{tabular}{|c|l|c|}
    \hline
    \textbf{เซต $A$} & \textbf{เพาเวอร์เซต $\mathcal{P}(A)$} & \textbf{จำนวนสมาชิก $|\mathcal{P}(A)|$} \\ \hline
    $\emptyset$ & $\{\emptyset\}$ & $2^0 = 1$ \\ \hline
    $\{a\}$ & $\{\emptyset, \{a\}\}$ & $2^1 = 2$ \\ \hline
    $\{1, 2\}$ & $\{\emptyset, \{1\}, \{2\}, \{1, 2\}\}$ & $2^2 = 4$ \\ \hline
    $\{1, 2, 3\}$ & $\{\emptyset, \{1\}, \{2\}, \{3\}, \{1,2\}, \{1,3\}, \{2,3\}, \{1,2,3\}\}$ & $2^3 = 8$ \\ \hline
  \end{tabular}
  \end{center}
  \vspace{0.2em}
  \begin{exampleblock}{ข้อสังเกตเชิงโครงสร้าง}
    เมื่อเพิ่มสมาชิกใหม่เข้ามา 1 ตัว เซตย่อยชุดเดิมจะยังคงอยู่ครบ ($2^n$ เซตเดิม)
    และจะเกิดเซตย่อยชุดใหม่ที่เติมสมาชิกตัวนั้นเข้าไปอีกเท่าตัว ($2^n$ เซตใหม่)
    ทำให้จำนวนสมาชิกรวมเพิ่มเป็น \textbf{2 เท่าเสมอ} ($2^n + 2^n = 2^{n+1}$)
  \end{exampleblock}
\end{frame}

\begin{frame}{ทำไม $|\mathcal{P}(A)| = 2^{|A|}$ ? (มุมมองคณิตศาสตร์)}
  \begin{columns}[T]
    \begin{column}{0.48\textwidth}
      \begin{block}{1. กฎการคูณ (Multiplication)}
        ในการสร้างเซตย่อย สมาชิกแต่ละตัวมีทางเลือกตัดสินใจ \textbf{2 ทาง}:
        \begin{itemize}
          \item \textbf{อยู่} ในเซตย่อย (1)
          \item \textbf{ไม่อยู่} ในเซตย่อย (0)
        \end{itemize}
        มีสมาชิก $n$ ตัว เป็นอิสระต่อกัน:
        \[ \underbrace{2 \times 2 \times \cdots \times 2}_{n \text{ ครั้ง}} = 2^n \text{ แบบ} \]
      \end{block}
    \end{column}
    \begin{column}{0.48\textwidth}
      \begin{block}{2. ทฤษฎีบททวินาม (Binomial)}
        นับแยกตามขนาดเซตย่อย:
        \begin{itemize}\small
          \item เซตย่อยขนาด $0$ มี $\binom{n}{0}$ เซต
          \item เซตย่อยขนาด $k$ มี $\binom{n}{k}$ เซต
          \item เซตย่อยขนาด $n$ มี $\binom{n}{n}$ เซต
        \end{itemize}
        รวมทุกขนาดโดยใช้ทฤษฎีบททวินาม:
        \[ |\mathcal{P}(A)| = \sum_{k=0}^n \binom{n}{k} = (1+1)^n = 2^n \]
      \end{block}
    \end{column}
  \end{columns}
\end{frame}

\begin{frame}{ทำไม $|\mathcal{P}(A)| = 2^{|A|}$ ? (มุมมองบิตมาสก์ในคอมพิวเตอร์)}
  แทนการเลือกสมาชิกเข้าเซตย่อยด้วยรหัสฐานสองความยาว $n$ บิต (\textbf{Bitmask})
  \begin{center}\small
  \begin{tabular}{|c|c|c|c|l|}
    \hline
    \textbf{บิตมาสก์} & สมาชิก $1$ & สมาชิก $2$ & สมาชิก $3$ & \textbf{เซตย่อยที่ได้} \\ \hline
    \texttt{000} & 0 & 0 & 0 & $\emptyset$ (เซตว่าง) \\ \hline
    \texttt{100} & 1 & 0 & 0 & $\{1\}$ \\ \hline
    \texttt{010} & 0 & 1 & 0 & $\{2\}$ \\ \hline
    \texttt{001} & 0 & 0 & 1 & $\{3\}$ \\ \hline
    \texttt{110} & 1 & 1 & 0 & $\{1, 2\}$ \\ \hline
    \texttt{101} & 1 & 0 & 1 & $\{1, 3\}$ \\ \hline
    \texttt{011} & 0 & 1 & 1 & $\{2, 3\}$ \\ \hline
    \texttt{111} & 1 & 1 & 1 & $\{1, 2, 3\}$ ($A$ ทั้งเซต) \\ \hline
  \end{tabular}
  \end{center}
  \vspace{0.2em}
  \footnotesize รหัสที่เป็นไปได้ทั้งหมดมีตั้งแต่ $00\ldots0$ ถึง $11\ldots1$ รวม $2^n$ สถานะพอดี \\
  นี่คือเหตุผลที่สัญกรณ์ $2^A$ สื่อความหมายว่า $|\mathcal{P}(A)| = 2^{|A|}$
\end{frame}

\begin{frame}{การเติบโตเชิงเลขชี้กำลัง (Exponential Growth) และผลต่อขั้นตอนวิธี}
  \begin{itemize}
    \item เพาเวอร์เซตเติบโตแบบ \textbf{Exponential ($2^n$)} ทุกครั้งที่ $|A|$ เพิ่ม 1 ขนาดจะเบิ้ลเป็นสองเท่า
  \end{itemize}
  \begin{center}\small
  \begin{tabular}{|c|c|l|}
    \hline
    \textbf{ขนาดเซต $n$} & \textbf{ขนาดเพาเวอร์เซต $2^n$} & \textbf{การเปรียบเทียบในทางปฏิบัติ} \\ \hline
    $n = 3$ & $2^3 = 8$ & คำนวณด้วยมือได้สบาย \\ \hline
    $n = 10$ & $2^{10} = 1,024$ & ระดับกิโลไบต์ คอมพิวเตอร์ประมวลผลทันที \\ \hline
    $n = 20$ & $2^{20} = 1,048,576$ & มากกว่า \textbf{1 ล้านเซตย่อย} \\ \hline
    $n = 30$ & $2^{30} = 1,073,741,824$ & มากกว่า \textbf{1 พันล้าน} เริ่มใช้เวลาหลายวินาที \\ \hline
    $n = 100$ & $2^{100} \approx 1.27 \times 10^{30}$ & มากกว่าจำนวนอะตอมในโลก! \\ \hline
  \end{tabular}
  \end{center}
  \vspace{0.2em}
  \begin{alertblock}{ผลต่อการเขียนโปรแกรมและขั้นตอนวิธี}
    ขั้นตอนวิธีแบบ Brute-force ที่ต้องวนลูปตรวจสอบทุกเซตย่อย (เช่น ปัญหา Subset Sum หรือ Knapsack) จะมีความซับซ้อน $O(2^n)$ ซึ่ง \textbf{ไม่สามารถใช้งานได้กับข้อมูลขนาดใหญ่}
  \end{alertblock}
\end{frame}

\begin{frame}{ทฤษฎีบทของคันทอร์ (Cantor's Theorem) และเพาเวอร์เซตซ้อน}
  \begin{block}{ทฤษฎีบทของคันทอร์ (Cantor's Theorem)}
    สำหรับทุกเซต $A$ (ไม่ว่าจะเป็นเซตจำกัดหรือเซตอนันต์) จะได้ว่า
    \[ |\mathcal{P}(A)| > |A| \]
    นั่นคือ \textbf{เพาเวอร์เซตของเซตใด ๆ มีขนาดมากกว่าเซตตั้งต้นเสมอ}
  \end{block}
  \vspace{0.2em}
  \textbf{ผลลัพธ์สำคัญ}: ไม่มีเซตอนันต์ใดที่ใหญ่ที่สุด เพราะเราสามารถสร้างเซตอนันต์ที่ใหญ่กว่าเดิมได้เสมอด้วยการหาเพาเวอร์เซต $\mathcal{P}(A), \mathcal{P}(\mathcal{P}(A)), \ldots$
  \vspace{0.2em}
  \begin{exampleblock}{ตัวอย่าง : เพาเวอร์เซตซ้อน (Nested Power Sets)}
    \begin{itemize}\small
      \item $\mathcal{P}(\emptyset) = \{\emptyset\}$ \quad $\Rightarrow\ |\mathcal{P}(\emptyset)| = 2^0 = 1$
      \item $\mathcal{P}(\mathcal{P}(\emptyset)) = \mathcal{P}(\{\emptyset\}) = \{\emptyset, \{\emptyset\}\}$ \quad $\Rightarrow\ |\mathcal{P}(\mathcal{P}(\emptyset))| = 2^1 = 2$
      \item $|\mathcal{P}(\mathcal{P}(\mathcal{P}(\emptyset)))| = 2^2 = 4$
      \item ถ้า $|A| = n$ จะได้ว่า $|\mathcal{P}(\mathcal{P}(A))| = 2^{|\mathcal{P}(A)|} = 2^{2^n}$
    \end{itemize}
  \end{exampleblock}
\end{frame}

\begin{frame}{เช็กความเข้าใจ : เพาเวอร์เซต}
  \begin{enumerate}
    \item จงหา $\mathcal{P}(\{a, b\})$
    \item จงหา $|\mathcal{P}(\{1, 2, 3, 4, 5, 6\})|$ โดยไม่ต้องแจกแจงสมาชิก
    \item ข้อความต่อไปนี้ จริง หรือ เท็จ:
          \begin{enumerate}
            \item[\textbf{(ก)}] $\{1\} \in \mathcal{P}(\{1, 2\})$
            \item[\textbf{(ข)}] $\{1\} \subseteq \mathcal{P}(\{1, 2\})$
            \item[\textbf{(ค)}] $\emptyset \in \mathcal{P}(A)$ และ $\emptyset \subseteq \mathcal{P}(A)$ สำหรับทุกเซต $A$
          \end{enumerate}
  \end{enumerate}
  \pause
  \begin{exampleblock}{เฉลยและเหตุผล}
    \footnotesize
    \textbf{1.} $\{\emptyset, \{a\}, \{b\}, \{a,b\}\}$ \\
    \textbf{2.} $2^6 = 64$ สมาชิก \\
    \textbf{3.}
    \textbf{(ก)} \TT\ (เพราะ $\{1\} \subseteq \{1,2\}$ จึงเป็นสมาชิกของเพาเวอร์เซต) \quad
    \textbf{(ข)} \FF\ (สมาชิกของ $\{1\}$ คือ $1$ ซึ่งไม่ใช่เซตย่อยของ $\{1,2\}$) \quad
    \textbf{(ค)} \TT\ ($\emptyset \subseteq A \Rightarrow \emptyset \in \mathcal{P}(A)$ และเซตว่างเป็นเซตย่อยของทุกเซตเสมอ)
  \end{exampleblock}
\end{frame}


%% =====================================================================
\section{ผลคูณคาร์ทีเซียน (Cartesian Product)}
%% =====================================================================

\begin{frame}{คู่อันดับ (Ordered Pair) และผลคูณคาร์ทีเซียน}
  \begin{block}{คู่อันดับ (Ordered Pair)}
    \textbf{คู่อันดับ} $(a, b)$ มีลำดับตำแหน่งชัดเจน: $a$ เป็นสมาชิกตัวหน้า และ $b$ เป็นสมาชิกตัวหลัง
    \begin{itemize}
      \item $(a, b) \neq (b, a)$ ยกเว้นกรณี $a = b$ (ต่างจากเซต $\{a, b\} = \{b, a\}$)
      \item \textbf{การเท่ากัน}: $(a, b) = (c, d) \iff a = c \text{ และ } b = d$
    \end{itemize}
  \end{block}
  \vspace{0.2em}
  \begin{block}{นิยาม : ผลคูณคาร์ทีเซียน (Cartesian Product)}
    ผลคูณคาร์ทีเซียนของเซต $A$ และ $B$ เขียนแทนด้วย $A \times B$ คือ
    \[ A \times B = \{(a, b) \mid a \in A \text{ และ } b \in B\} \]
  \end{block}
  \vspace{0.2em}
  \footnotesize ตั้งชื่อตาม \textbf{Ren\'e Descartes} (เรอเน เดการ์ต) ผู้คิดค้นระบบพิกัดฉาก \\
  $\mathbb{R} \times \mathbb{R} = \mathbb{R}^2$ คือระนาบพิกัด 2 มิติ $(x, y)$ ที่เราคุ้นเคย
\end{frame}

\begin{frame}{ตัวอย่างและการเขียนตารางคาร์ทีเซียน (Cartesian Grid)}
  กำหนด $A = \{1, 2\}$ และ $B = \{a, b, c\}$
  \begin{center}
  \renewcommand{\arraystretch}{1.2}
  \begin{tabular}{|c|c|c|c|}
    \hline
    $A \times B$ & $a$ & $b$ & $c$ \\ \hline
    $1$ & $(1,a)$ & $(1,b)$ & $(1,c)$ \\ \hline
    $2$ & $(2,a)$ & $(2,b)$ & $(2,c)$ \\ \hline
  \end{tabular}
  \end{center}
  \begin{itemize}\small
    \item $A \times B = \{(1,a), (1,b), (1,c), (2,a), (2,b), (2,c)\}$
    \item $B \times A = \{(a,1), (a,2), (b,1), (b,2), (c,1), (c,2)\}$
  \end{itemize}
  \vspace{0.1em}
  \begin{alertblock}{ข้อสังเกตสำคัญ}
    \begin{itemize}\small
      \item $A \times B \neq B \times A$ (ผลคูณคาร์ทีเซียน \textbf{ไม่มีสมบัติการสลับที่})
      \item สำหรับเซตใด ๆ $A$: \quad $A \times \emptyset = \emptyset \times A = \emptyset$ (เพราะจับคู่ไม่ได้)
    \end{itemize}
  \end{alertblock}
\end{frame}

\begin{frame}{จำนวนสมาชิกของผลคูณคาร์ทีเซียน และการขยาย $k$ มิติ}
  \begin{block}{สูตรจำนวนสมาชิก}
    ถ้า $|A| = m$ และ $|B| = n$ แล้ว
    \[ |A \times B| = |A| \cdot |B| = m \cdot n \]
  \end{block}
  \vspace{0.2em}
  \textbf{การขยายสู่ผลคูณคาร์ทีเซียนหลายเซต ($k$-tuple)}:
  \[ A_1 \times A_2 \times \cdots \times A_k = \{(a_1, a_2, \ldots, a_k) \mid a_i \in A_i\} \]
  \[ |A_1 \times A_2 \times \cdots \times A_k| = |A_1| \cdot |A_2| \cdots |A_k| \]
  \vspace{0.2em}
  \begin{exampleblock}{ตัวอย่างในระบบคอมพิวเตอร์}
    \begin{itemize}\small
      \item $\{0, 1\}^3 = \{0, 1\} \times \{0, 1\} \times \{0, 1\}$ มีสมาชิก $2 \times 2 \times 2 = 8$ ตัว
            ได้แก่ $(0,0,0), (0,0,1), \ldots, (1,1,1)$ (พิกัดบิต 3 มิติ)
      \item จานสี RGB: $\text{Red} \times \text{Green} \times \text{Blue} = \{0..255\}^3 \Rightarrow 256^3 = 16,777,216$ สี
    \end{itemize}
  \end{exampleblock}
\end{frame}

\begin{frame}{สมบัติการกระจายของผลคูณคาร์ทีเซียน}
  \begin{block}{ทฤษฎีบท : สมบัติการกระจาย (Distributive Properties)}
    สำหรับเซต $A, B, C$ ใด ๆ จะได้ว่า
    \begin{enumerate}
      \item $A \times (B \cup C) = (A \times B) \cup (A \times C)$
      \item $A \times (B \cap C) = (A \times B) \cap (A \times C)$
      \item $(A \cup B) \times C = (A \times C) \cup (B \times C)$
      \item $(A \cap B) \times C = (A \times C) \cap (B \times C)$
    \end{enumerate}
  \end{block}
  \vspace{0.2em}
  \begin{exampleblock}{ความคล้ายคลึงกับพีชคณิตตัวเลข}
    สังเกตว่าผลคูณคาร์ทีเซียน ($\times$) กระจายเหนือยูเนียน ($\cup$) และอินเตอร์เซกชัน ($\cap$)
    ในลักษณะเดียวกับการคูณตัวเลขที่กระจายเหนือการบวก: $a \cdot (b + c) = a \cdot b + a \cdot c$
  \end{exampleblock}
\end{frame}

\begin{frame}{ผลคูณคาร์ทีเซียนกับระบบฐานข้อมูลเชิงสัมพันธ์ (RDBMS)}
  \begin{itemize}
    \item ให้ $S = \{\text{นักศึกษา}\}$ และ $C = \{\text{รายวิชา}\}$
    \item $S \times C$ คือการจับคู่นักศึกษาทุกคนกับทุกวิชาที่เป็นไปได้
    \item \textbf{ตารางลงทะเบียน (Enrollment Table)} คือเซตย่อยของผลคูณคาร์ทีเซียน:
          \[ E \subseteq S \times C \]
    \item ในภาษา SQL คำสั่ง \texttt{CROSS JOIN} คือการหา $A \times B$ โดยตรง:
          \begin{center}
          \texttt{SELECT * FROM Students CROSS JOIN Courses;}
          \end{center}
    \item แต่ละแถว (Row/Record) ในตารางฐานข้อมูลคือ \textbf{Tuple} หนึ่งตัวในผลคูณคาร์ทีเซียน
  \end{itemize}
  \vspace{0.2em}
  \begin{exampleblock}{สะพานเชื่อมสู่บทที่ 3}
    ในบทที่ 3 เราจะนิยาม \textbf{``ความสัมพันธ์ (Relation)''} และ \textbf{``ฟังก์ชัน (Function)''}
    ว่าเป็น \textbf{เซตย่อยของผลคูณคาร์ทีเซียน $A \times B$} อย่างเป็นทางการ
  \end{exampleblock}
\end{frame}

\begin{frame}{เช็กความเข้าใจ : ผลคูณคาร์ทีเซียน}
  กำหนด $A = \{1, 2\}$, $B = \{x, y\}$ และ $C = \{y, z\}$
  \begin{enumerate}
    \item จงหา $A \times B$ และ $B \times A$ พร้อมเปรียบเทียบว่าเท่ากันหรือไม่
    \item ถ้า $(2x - 1, \ 5) = (7, \ y + 2)$ จงหาค่าของ $x$ และ $y$
    \item ถ้า $|A| = 4, |B| = 3, |C| = 2$ จงหา $|A \times B \times C|$ และ $|\mathcal{P}(A \times C)|$
    \item จงหา $A \times (B \cap C)$
  \end{enumerate}
  \pause
  \begin{exampleblock}{เฉลย}
    \footnotesize
    \textbf{1.} $A \times B = \{(1,x),(1,y),(2,x),(2,y)\}$, \ $B \times A = \{(x,1),(x,2),(y,1),(y,2)\}$ (ไม่เท่ากัน) \\
    \textbf{2.} เทียบคู่อันดับ: $2x - 1 = 7 \Rightarrow x = 4$ และ $y + 2 = 5 \Rightarrow y = 3$ \\
    \textbf{3.} $|A \times B \times C| = 4 \times 3 \times 2 = 24$ \quad และ \quad $|\mathcal{P}(A \times C)| = 2^{|A \times C|} = 2^{4 \times 2} = 2^8 = 256$ \\
    \textbf{4.} $B \cap C = \{y\} \ \Rightarrow\ A \times (B \cap C) = \{1,2\} \times \{y\} = \{(1,y), (2,y)\}$
  \end{exampleblock}
\end{frame}


%% =====================================================================
\section{จำนวนสมาชิกและหลักการรวม-การแยก}
%% =====================================================================

\begin{frame}{จำนวนสมาชิกของเซต (Cardinality)}
  \begin{block}{นิยาม : คาร์ดินัลลิตี้ (Cardinality)}
    \textbf{จำนวนสมาชิกของเซต (Cardinality)} ของเซต $A$ เขียนแทนด้วย $|A|$ หรือ $n(A)$
    คือ \textbf{จำนวนสมาชิกที่แตกต่างกันทั้งหมด} ของเซต $A$
  \end{block}
  \vspace{0.2em}
  \begin{itemize}
    \item \textbf{เซตจำกัด (Finite Set)}: นับจำนวนสมาชิกได้เป็นจำนวนเต็มบวกหรือศูนย์ เช่น $|\{a,e,i,o,u\}| = 5, |\emptyset| = 0$
    \item \textbf{เซตอนันต์ (Infinite Set)}: สมาชิกนับไม่ถ้วน เปรียบเทียบขนาดด้วยการจับคู่ 1-1 (เช่น $|\mathbb{N}|, |\mathbb{R}|$)
  \end{itemize}
  \vspace{0.2em}
  \begin{center}\small
  \begin{tabular}{|l|c|l|}
    \hline
    \textbf{โครงสร้างเซต} & \textbf{สูตรจำนวนสมาชิก} & \textbf{เงื่อนไข/คำอธิบาย} \\ \hline
    เพาเวอร์เซต & $|\mathcal{P}(A)| = 2^{|A|}$ & สมาชิกแต่ละตัวมี 2 ทางเลือก \\ \hline
    ผลคูณคาร์ทีเซียน & $|A \times B| = |A| \cdot |B|$ & กฎการคูณการจับคู่ \\ \hline
    เซตไม่มีสมาชิกร่วม & $|A \cup B| = |A| + |B|$ & เมื่อ $A \cap B = \emptyset$ (Disjoint Sets) \\ \hline
  \end{tabular}
  \end{center}
\end{frame}

\begin{frame}{หลักการรวม-การแยกสำหรับ 2 เซต (Inclusion--Exclusion Principle)}
  \begin{alertblock}{ทำไม $|A \cup B| \neq |A| + |B|$ โดยทั่วไป ?}
    ถ้านับ $|A| + |B|$ สมาชิกที่อยู่ในส่วนซ้อนทับ ($A \cap B$) จะถูก \textbf{นับซ้ำ 2 ครั้ง} จึงต้องลบออก 1 ครั้ง
  \end{alertblock}
  \begin{columns}[T]
    \begin{column}{0.54\textwidth}
      \begin{block}{สูตรหลัก (2 เซต)}
        \[ |A \cup B| = |A| + |B| - |A \cap B| \]
      \end{block}
      \begin{block}{สูตรอนุพันธ์ที่ใช้บ่อย}
        \begin{itemize}\small
          \item $|A \cap B| = |A| + |B| - |A \cup B|$
          \item เฉพาะ $A$: $|A - B| = |A| - |A \cap B|$
          \item เฉพาะ $B$: $|B - A| = |B| - |A \cap B|$
          \item คอมพลีเมนต์: $|A^c| = |U| - |A|$
          \item นอกทั้งคู่: $|(A \cup B)^c| = |U| - |A \cup B|$
        \end{itemize}
      \end{block}
    \end{column}
    \begin{column}{0.44\textwidth}
      \begin{center}
      \begin{tikzpicture}[scale=0.8,every node/.style={font=\small}]
        \draw[thick] (-2.2,-1.4) rectangle (2.2,1.5);
        \node at (1.9,1.2) {$U$};
        \draw[fill=lightnavy,fill opacity=0.7] (-0.7,0) circle (1.1);
        \draw[fill=mint,fill opacity=0.7] (0.7,0) circle (1.1);
        \draw[thick] (-0.7,0) circle (1.1);
        \draw[thick] (0.7,0) circle (1.1);
        \node at (-1.2,0) {\footnotesize $|A-B|$};
        \node at (0,0) {\footnotesize $|A\cap B|$};
        \node at (1.2,0) {\footnotesize $|B-A|$};
        \node at (-1.4,1.2) {$A$};
        \node at (1.4,1.2) {$B$};
      \end{tikzpicture}\\
      \footnotesize ส่วนซ้ำตรงกลางนับซ้ำ ต้องหักออกหนึ่งครั้ง
      \end{center}
    \end{column}
  \end{columns}
\end{frame}

\begin{frame}{ตัวอย่าง : การแก้โจทย์ 2 เซตอย่างเป็นระบบ}
  \begin{block}{โจทย์}
    ในการสำรวจนักศึกษา 100 คน พบว่า:
    มีผู้ใช้ภาษา Python 65 คน, ใช้ภาษา Java 45 คน, และใช้ \textbf{ทั้งสองภาษา} 25 คน
    \begin{enumerate}
      \item มีนักศึกษากี่คนที่ใช้ \textbf{อย่างน้อยหนึ่งภาษา}?
      \item มีนักศึกษากี่คนที่ \textbf{ไม่ใช้ทั้งสองภาษาเลย}?
      \item มีนักศึกษากี่คนที่ใช้ \textbf{ภาษา Python เพียงภาษาเดียว}?
    \end{enumerate}
  \end{block}
  \pause
  \textbf{วิธีทำ} \quad ให้ $|U| = 100$, $P$ = กลุ่มที่ใช้ Python ($|P| = 65$), $J$ = กลุ่มที่ใช้ Java ($|J| = 45$), $|P \cap J| = 25$
  \begin{align*}
    \text{1.}\quad |P \cup J| &= |P| + |J| - |P \cap J| = 65 + 45 - 25 = \mathbf{85} \text{ คน} \\
    \text{2.}\quad |(P \cup J)^c| &= |U| - |P \cup J| = 100 - 85 = \mathbf{15} \text{ คน} \\
    \text{3.}\quad |P - J| &= |P| - |P \cap J| = 65 - 25 = \mathbf{40} \text{ คน}
  \end{align*}
  \begin{exampleblock}{ตรวจคำตอบด้วยแผนภาพเวนน์}
    ช่องซ้าย ($40$) + ช่องกลาง ($25$) + ช่องขวา ($20$) + ช่องนอก ($15$) $= 100$ คน พอดี!
  \end{exampleblock}
\end{frame}

\begin{frame}{หลักการรวม-การแยกสำหรับ 3 เซต (Inclusion--Exclusion for 3 Sets)}
  \begin{block}{ทฤษฎีบท : สูตร 3 เซต}
    สำหรับเซตจำกัด $A, B, C$ ใด ๆ
    \[
      |A \cup B \cup C| = \underbrace{(|A| + |B| + |C|)}_{\text{บวกเซตเดี่ยว}} - \underbrace{(|A \cap B| + |A \cap C| + |B \cap C|)}_{\text{ลบส่วนซ้อน 2 เซต}} + \underbrace{|A \cap B \cap C|}_{\text{บวกส่วนซ้อน 3 เซตกลับคืน}}
    \]
  \end{block}
  \vspace{0.2em}
  \textbf{ทำไมต้องบวกกลับในขั้นตอนสุดท้าย? (Logic Trace)}
  \begin{enumerate}\small
    \item เมื่อบวก $|A| + |B| + |C|$ : ส่วน $A \cap B \cap C$ ถูกบวก \textbf{3 ครั้ง}
    \item เมื่อลบส่วนซ้อนคู่ 3 คู่ : ส่วน $A \cap B \cap C$ ถูกลบออก \textbf{3 ครั้ง} ($+3 - 3 = 0$) ทำให้ส่วนกลาง \textbf{หายไปหมด!}
    \item จึงจำเป็นต้อง \textbf{บวก $|A \cap B \cap C|$ กลับเข้าไป 1 ครั้ง} เพื่อให้นับครบพอดี
  \end{enumerate}
\end{frame}

\begin{frame}{ตัวอย่าง : โจทย์ 3 เซต (วิชาเรียน)}
  \begin{block}{โจทย์}
    ห้องเรียน 40 คน: เรียนคณิต ($M$) 22 คน, ฟิสิกส์ ($P$) 20 คน, เคมี ($C$) 15 คน \\
    เรียนคณิตและฟิสิกส์ 10 คน, คณิตและเคมี 7 คน, ฟิสิกส์และเคมี 6 คน, และเรียนทั้งสามวิชา 3 คน \\
    จงหาจำนวนผู้ที่เรียนอย่างน้อยหนึ่งวิชา และจำนวนผู้ที่ไม่เรียนวิชาใดเลย
  \end{block}
  \pause
  \textbf{วิธีทำ} \quad แทนค่าลงในสูตร 3 เซตโดยตรง:
  \begin{align*}
    |M \cup P \cup C| &= |M| + |P| + |C| - |M \cap P| - |M \cap C| - |P \cap C| + |M \cap P \cap C| \\
                      &= 22 + 20 + 15 - 10 - 7 - 6 + 3 = \mathbf{37} \text{ คน}
  \end{align*}
  \begin{itemize}
    \item ผู้ที่เรียนอย่างน้อยหนึ่งวิชา $= \mathbf{37}$ คน
    \item ผู้ที่ไม่เรียนวิชาใดเลย $= |U| - |M \cup P \cup C| = 40 - 37 = \mathbf{3}$ คน
  \end{itemize}
\end{frame}

\begin{frame}{การตรวจสอบความสอดคล้องของข้อมูล (Data Validation)}
  หลักการรวม-การแยกไม่ได้มีไว้เพียงคำนวณตัวเลข แต่เป็น \textbf{เครื่องมือตรวจสอบความสอดคล้องของข้อมูล (Data Consistency)}
  \begin{block}{ตัวอย่างการจับผิดข้อมูล (Data Inconsistency)}
    ผู้สำรวจรายงานว่า ในห้องเรียน 20 คน:
    เรียนคณิต 12 คน, ฟิสิกส์ 15 คน, เคมี 10 คน, คณิต$\cap$ฟิสิกส์ 6 คน, คณิต$\cap$เคมี 5 คน, ฟิสิกส์$\cap$เคมี 4 คน, และเรียนทั้งสามวิชา 2 คน
  \end{block}
  \pause
  คำนวณตามสูตร:
  \[ |M \cup P \cup C| = 12 + 15 + 10 - 6 - 5 - 4 + 2 = \mathbf{24} \text{ คน} \]
  \begin{alertblock}{ผลการตรวจสอบ : ข้อมูลขัดแย้ง!}
    ผลลัพธ์ได้ $24$ คน แต่ห้องเรียนมีนักเรียนทั้งหมดเพียง $20$ คน ($|M \cup P \cup C| > |U|$) \\
    $\Rightarrow$ \textbf{ข้อมูลรายงานชุดนี้เป็นเท็จ / ผิดพลาด} ต้องปฏิเสธข้อมูลและส่งตรวจสอบใหม่
  \end{alertblock}
\end{frame}


%% =====================================================================
\section{เซตในวิทยาการคอมพิวเตอร์}
%% =====================================================================

\begin{frame}{การแทนเซตด้วยเวกเตอร์บิต (Bit Vector Representation)}
  ในเอกภพสัมพัทธ์จำกัด $U = \{1, 2, \ldots, n\}$ เราแทนเซตด้วย \textbf{บิต $n$ ตำแหน่ง}:
  บิตที่ $i$ มีค่าเป็น \texttt{1} เมื่อ $i \in A$ และเป็น \texttt{0} เมื่อ $i \notin A$
  \begin{block}{ตัวอย่าง: $U = \{1, 2, 3, 4, 5, 6, 7, 8, 9\}$ และ $A = \{2, 5, 7\}$}
    \centering
    \begin{tabular}{r|ccccccccc}
      ตำแหน่งใน $U$ & 1 & 2 & 3 & 4 & 5 & 6 & 7 & 8 & 9 \\ \hline
      บิตของ $A$ & 0 & \textbf{\textcolor{navy}{1}} & 0 & 0 & \textbf{\textcolor{navy}{1}} & 0 & \textbf{\textcolor{navy}{1}} & 0 & 0 \\
    \end{tabular}
    \par\vspace{0.2em} เวกเตอร์บิตของ $A = \texttt{010010100}_2$
  \end{block}
  \vspace{0.2em}
  การดำเนินการเซตแปลงเป็นการดำเนินการบิตระดับฮาร์ดแวร์ (Bitwise Operations) ที่เร็วที่สุด:
  \begin{center}\small
  \begin{tabular}{|l|c|l|c|}
    \hline
    \textbf{การดำเนินการเซต} & \textbf{ตัวดำเนินการบิต} & \textbf{ไวยากรณ์ภาษา C/Java/Python} & \textbf{ความเร็ว} \\ \hline
    ยูเนียน $A \cup B$ & OR & \texttt{A | B} & $O(1)$ ต่อ word \\ \hline
    อินเตอร์เซกชัน $A \cap B$ & AND & \texttt{A \& B} & $O(1)$ ต่อ word \\ \hline
    ผลต่าง $A - B$ & AND NOT & \texttt{A \& (\textasciitilde B)} & $O(1)$ ต่อ word \\ \hline
    คอมพลีเมนต์ $A^c$ & NOT & \texttt{\textasciitilde A} & $O(1)$ ต่อ word \\ \hline
  \end{tabular}
  \end{center}
\end{frame}

\begin{frame}{เซตในภาษา SQL และระบบฐานข้อมูล}
  ทฤษฎีเซตคือรากฐานของโมเดลฐานข้อมูลเชิงสัมพันธ์ (Relational Model)
  \begin{center}\small
  \begin{tabular}{|l|l|l|}
    \hline
    \textbf{ทฤษฎีเซต} & \textbf{คำสั่ง SQL} & \textbf{ความหมายในฐานข้อมูล} \\ \hline
    $A \cup B$ & \texttt{SELECT ... UNION ...} & รวมแถวจากสองตาราง (ตัดตัวซ้ำ) \\ \hline
    $A \cap B$ & \texttt{SELECT ... INTERSECT ...} & เอาเฉพาะแถวที่ปรากฏในทั้งสองตาราง \\ \hline
    $A - B$ & \texttt{SELECT ... EXCEPT ...} & เอาแถวในตารางแรกที่ไม่ปรากฏในตารางสอง \\ \hline
    $A \times B$ & \texttt{SELECT ... CROSS JOIN ...} & จับคู่ทุกแถวของ $A$ กับทุกแถวของ $B$ \\ \hline
    $\{x \in A \mid P(x)\}$ & \texttt{SELECT ... WHERE condition} & กรองสมาชิกย่อยตามเงื่อนไขเพรดิเคต \\ \hline
  \end{tabular}
  \end{center}
  \vspace{0.2em}
  \begin{exampleblock}{ตัวอย่างจริง}
    หานักศึกษาที่ลงทะเบียนวิชาคณิตศาสตร์แต่ไม่ลงทะเบียนวิชาคอมพิวเตอร์:
    \begin{center}
      \texttt{SELECT sid FROM enroll\_math EXCEPT SELECT sid FROM enroll\_cs;}
    \end{center}
  \end{exampleblock}
\end{frame}

\begin{frame}{ตารางแฮช (Hash Set) : ประสิทธิภาพ $O(1)$}
  ในภาษาโปรแกรมสมัยใหม่ เซตถูกพัฒนาด้วย \textbf{ตารางแฮช (Hash Table)}
  โดยใช้ฟังก์ชันแฮช $h(x)$ แปลงสมาชิกเป็นตำแหน่งดัชนีในหน่วยความจำโดยตรง
  \begin{center}\small
  \begin{tabular}{|l|c|c|}
    \hline
    \textbf{การดำเนินการ} & \textbf{Hash Set (เวลาเฉลี่ย)} & \textbf{Array / List ทั่วไป} \\ \hline
    ทดสอบสมาชิก $x \in A$ & $\mathbf{O(1)}$ & $O(n)$ \\ \hline
    เพิ่มสมาชิก (Insert) & $\mathbf{O(1)}$ & $O(1)$ หรือ $O(n)$ \\ \hline
    ลบสมาชิก (Delete) & $\mathbf{O(1)}$ & $O(n)$ \\ \hline
    หายูเนียน $A \cup B$ & $\mathbf{O(|A| + |B|)}$ & $O(|A| \cdot |B|)$ \\ \hline
    หาอินเตอร์เซกชัน $A \cap B$ & $\mathbf{O(\min(|A|, |B|))}$ & $O(|A| \cdot |B|)$ \\ \hline
  \end{tabular}
  \end{center}
  \vspace{0.2em}
  \begin{itemize}
    \item ภาษาที่รองรับ: Python (\texttt{set}), Java (\texttt{HashSet}), C++ (\texttt{std::unordered\_set})
    \item การเลือกโครงสร้างข้อมูลเซตที่ถูกต้องช่วยเพิ่มความเร็วการทำงานของระบบได้มหาศาล
  \end{itemize}
\end{frame}

\begin{frame}{เซตนับได้และเซตนับไม่ได้ (Countable vs Uncountable Sets)}
  \begin{columns}[T]
    \begin{column}{0.48\textwidth}
      \begin{block}{โรงแรมของฮิลเบิร์ต (Hilbert's Hotel)}
        โรงแรมมีห้อง $1, 2, 3, \ldots$ อนันต์ห้อง \textbf{และเต็มทุกห้อง} \\
        เมื่อมีแขกใหม่มา 1 คน: ให้แขกห้อง $n$ ย้ายไปห้อง $n+1$ ทำให้ห้อง 1 ว่างลง! \\
        $\Rightarrow$ เซตอนันต์สามารถรับสมาชิกเพิ่มได้โดยขนาดไม่เปลี่ยน
      \end{block}
    \end{column}
    \begin{column}{0.50\textwidth}
      \begin{block}{ระดับของอนันต์ (Infinity Hierarchy)}
        \begin{itemize}\small
          \item \textbf{เซตนับได้ (Countable)}: จับคู่ 1-1 กับ $\mathbb{N}$ ได้
                \[ |\mathbb{N}| = |\mathbb{Z}| = |\mathbb{Q}| = \aleph_0 \]
          \item \textbf{เซตนับไม่ได้ (Uncountable)}: เช่น จำนวนจริง $\mathbb{R}$ (พิสูจน์โดย Cantor's Diagonal Argument)
                \[ |\mathbb{N}| < |\mathbb{R}| < |\mathcal{P}(\mathbb{R})| \]
        \end{itemize}
      \end{block}
    \end{column}
  \end{columns}
  \vspace{0.2em}
  \begin{exampleblock}{ข้อสรุป}
    อนันต์มีหลายขนาด และเพาเวอร์เซตสร้างอนันต์ขั้นที่สูงขึ้นไปได้ไม่รู้จบ
  \end{exampleblock}
\end{frame}


%% =====================================================================
\section{กิจกรรมในชั้นเรียน}
%% =====================================================================

\begin{frame}{กิจกรรมที่ 1 : เพาเวอร์เซต \eng{6 นาที}}
  \begin{enumerate}
    \item กำหนด $S = \{x, y, z\}$ จงแจกแจงสมาชิกทุกตัวของ $\mathcal{P}(S)$
    \item จำนวนสมาชิก $|\mathcal{P}(S)|$ มีกี่ตัว และตรงกับสูตร $2^n$ หรือไม่
    \item พิจารณาข้อความต่อไปนี้ว่า \textbf{จริง} หรือ \textbf{เท็จ}:
          \begin{itemize}
            \item[\textbf{(ก)}] $\emptyset \subseteq \mathcal{P}(S)$
            \item[\textbf{(ข)}] $\emptyset \in \mathcal{P}(S)$
            \item[\textbf{(ค)}] $\{x, y\} \subseteq \mathcal{P}(S)$
            \item[\textbf{(ง)}] $\{\{x, y\}\} \subseteq \mathcal{P}(S)$
          \end{itemize}
  \end{enumerate}
  \pause
  \begin{exampleblock}{เฉลยกิจกรรมที่ 1}
    \footnotesize
    \textbf{1.} $\mathcal{P}(S) = \{\emptyset, \{x\}, \{y\}, \{z\}, \{x,y\}, \{x,z\}, \{y,z\}, \{x,y,z\}\}$ \\
    \textbf{2.} มี 8 ตัว $= 2^3$ ตรงตามสูตรพอดี \\
    \textbf{3.} \textbf{(ก)} \TT\ ($\emptyset$ เป็นสับเซตของทุกเซต) \quad
    \textbf{(ข)} \TT\ ($\emptyset \subseteq S \Rightarrow \emptyset \in \mathcal{P}(S)$) \quad
    \textbf{(ค)} \FF\ ($x \notin \mathcal{P}(S)$ เพราะสมาชิกของเพาเวอร์เซตต้องเป็นเซต) \quad
    \textbf{(ง)} \TT\ (เพราะ $\{x,y\} \in \mathcal{P}(S)$)
  \end{exampleblock}
\end{frame}

\begin{frame}{กิจกรรมที่ 2 : ผลคูณคาร์ทีเซียน \eng{6 นาที}}
  กำหนด $A = \{1, 2\}$ และ $B = \{a, b, c\}$
  \begin{enumerate}
    \item จงแจกแจงสมาชิกของ $A \times B$ และวาดตารางคาร์ทีเซียน
    \item $|A \times B|$ และ $|\mathcal{P}(A \times B)|$ มีค่าเท่าใด
    \item $A \times B = B \times A$ หรือไม่ จงอธิบายเหตุผล
    \item จงหา $|\{1, 2\} \times \{3, 4\} \times \{5, 6\}|$
  \end{enumerate}
  \pause
  \begin{exampleblock}{เฉลยกิจกรรมที่ 2}
    \footnotesize
    \textbf{1.} $A \times B = \{(1,a), (1,b), (1,c), (2,a), (2,b), (2,c)\}$ \\
    \textbf{2.} $|A \times B| = 2 \times 3 = 6$ \quad และ \quad $|\mathcal{P}(A \times B)| = 2^6 = 64$ \\
    \textbf{3.} ไม่เท่ากัน เพราะ $(1,a) \in A \times B$ แต่ $(1,a) \notin B \times A$ เนื่องจากคู่อันดับมีลำดับหน้า-หลัง \\
    \textbf{4.} $2 \times 2 \times 2 = 8$ สมาชิก
  \end{exampleblock}
\end{frame}

\begin{frame}{กิจกรรมที่ 3 : การนับสมาชิกและตรวจความสอดคล้อง \eng{8 นาที}}
  \begin{block}{โจทย์กิจกรรม}
    ชมรมคอมพิวเตอร์มีสมาชิก 50 คน พบว่า:
    ชอบเขียนโปรแกรม ($P$) 30 คน, ชอบระบบเครือข่าย ($N$) 25 คน, และชอบทั้งสองอย่าง 12 คน
    \begin{enumerate}
      \item มีกี่คนที่ชอบอย่างน้อยหนึ่งอย่าง?
      \item มีกี่คนที่ไม่ชอบทั้งสองอย่างนี้เลย?
      \item มีกี่คนที่ชอบเขียนโปรแกรมเพียงอย่างเดียว?
      \item ถ้ามีผู้รายงานเพิ่มว่ามีคนชอบฐานข้อมูล 20 คน โดยชอบทั้ง 3 ด้าน 5 คน และชอบเครือข่ายกับฐานข้อมูล 15 คน ข้อมูลนี้เป็นไปได้หรือไม่
    \end{enumerate}
  \end{block}
  \pause
  \begin{exampleblock}{เฉลยกิจกรรมที่ 3}
    \footnotesize
    \textbf{1.} $|P \cup N| = 30 + 25 - 12 = \mathbf{43}$ คน \quad
    \textbf{2.} $50 - 43 = \mathbf{7}$ คน \quad
    \textbf{3.} $|P - N| = 30 - 12 = \mathbf{18}$ คน \\
    \textbf{4.} ตรวจสอบ: $|N \cap D| = 15 \leq |N| = 25$ และ $15 \leq |D| = 20$ (เป็นไปได้ในเบื้องต้น แต่ต้องตรวจทั้งระบบว่าไม่เกิน $|U|=50$)
  \end{exampleblock}
\end{frame}

\begin{frame}{เกมท้ายคาบ : ถอดรหัสลับบิต (Bit Vector Cipher)}
  ใช้เอกภพสัมพัทธ์ $U = \{1, 2, 3, 4, 5, 6, 7, 8\}$ (บิตซ้ายสุด = สมาชิกตัวที่ 1)
  \begin{enumerate}
    \item[\textbf{รอบ 1}] เวกเตอร์บิต \texttt{10100110} แทนเซตใด? \eng{1 แต้ม}
    \pause
    \item[] \hfill \textcolor{navy}{\textbf{เฉลย:}} $\{1, 3, 6, 7\}$
    \pause
    \item[\textbf{รอบ 2}] กำหนด $A = \texttt{11001100}$ และ $B = \texttt{01101000}$ \\
          จงหา $A \cap B$ (AND บิตต่อบิต) แล้วแปลงผลลัพธ์เป็นเซต \eng{2 แต้ม}
    \pause
    \item[] \hfill \textcolor{navy}{\textbf{เฉลย:}} $\texttt{11001100} \ \& \ \texttt{01101000} = \texttt{01001000} \Rightarrow \{2, 5\}$
    \pause
    \item[\textbf{รอบ 3}] กำหนดรหัสตัวอักษร: $1{=}\text{S}, 2{=}\text{E}, 3{=}\text{T}, 4{=}\text{M}, 5{=}\text{A}, 6{=}\text{T}, 7{=}\text{H}$ \\
          ถ้า $A = \texttt{1010100}, B = \texttt{0110001}, C = \texttt{0000101}$ \\
          จงหาคำลับจาก $(A \cup B) - C$ \eng{3 แต้ม}
    \pause
    \item[] \hfill \textcolor{navy}{\textbf{เฉลย:}} $(A \cup B) - C = \texttt{1110000} \Rightarrow \{1, 2, 3\} \Rightarrow \mathbf{\text{``SET''}}$
  \end{enumerate}
\end{frame}


%% ====================== สรุป + แบบฝึกหัด ======================
\begin{frame}{สรุปภาพรวมสัปดาห์ที่ 6 และบทที่ 2 ทฤษฎีเซต}
  \begin{center}\small
  \begin{tabular}{|l|l|l|}
    \hline
    \textbf{หัวข้อ} & \textbf{สูตร/นิยามสำคัญ} & \textbf{การประยุกต์ในคอมพิวเตอร์} \\ \hline
    \textbf{เพาเวอร์เซต} & $\mathcal{P}(A) = \{B \mid B \subseteq A\}, \ |\mathcal{P}(A)| = 2^{|A|}$ & บิตมาสก์, สถานะเซตย่อย, Exponential $O(2^n)$ \\ \hline
    \textbf{ผลคูณคาร์ทีเซียน} & $A \times B = \{(a,b) \mid a \in A, b \in B\}, \ |A \times B| = |A|\cdot|B|$ & ตารางฐานข้อมูล, \texttt{CROSS JOIN}, ฟังก์ชัน \\ \hline
    \textbf{รวม-การแยก (2 เซต)} & $|A \cup B| = |A| + |B| - |A \cap B|$ & การนับข้อมูลซ้ำ, Query Optimizer \\ \hline
    \textbf{รวม-การแยก (3 เซต)} & $|A \cup B \cup C| = \sum |A| - \sum |A \cap B| + |A \cap B \cap C|$ & ตรวจสอบความถูกต้องของข้อมูล (Validation) \\ \hline
    \textbf{โครงสร้างข้อมูล} & Bit Vector, Hash Table, SQL & Bitwise $O(1)$, Hash Set $O(1)$, Set Ops \\ \hline
  \end{tabular}
  \end{center}
  \vspace{0.2em}
  \begin{exampleblock}{บทที่ 2 ทฤษฎีเซต ใน 3 บรรทัด}
    \begin{itemize}\small
      \item สัปดาห์ 4: นิยามเซต สมาชิก เซตย่อย เอกภพสัมพัทธ์ เซตว่าง
      \item สัปดาห์ 5: ยูเนียน อินเตอร์เซกชัน ผลต่าง คอมพลีเมนต์ และกฎเดอมอร์แกน
      \item สัปดาห์ 6: เพาเวอร์เซต ผลคูณคาร์ทีเซียน หลักการนับ และเซตในระบบคอมพิวเตอร์
    \end{itemize}
  \end{exampleblock}
\end{frame}

\begin{frame}{แบบฝึกหัดมอบหมาย (สัปดาห์ที่ 6)}
  ทำลงสมุด ส่งต้นคาบสัปดาห์หน้า (ข้อละ 2 คะแนน รวม 10 คะแนน)
  \begin{enumerate}\small
    \item กำหนด $A = \{1, 2, 3\}$ จงหา $\mathcal{P}(A)$ และหาค่าของ $|\mathcal{P}(\mathcal{P}(\emptyset))|$
    \item กำหนด $A = \{a, b\}$ และ $B = \{1, 2, 3\}$ จงหา $A \times B$, $B \times A$ และ $|\mathcal{P}(A \times B)|$
    \item ถ้า $|A| = 3, |B| = 4, |C| = 2$ จงหา $|A \times B \times C|$ และ $|(A \times B) \cup (A \times C)|$ เมื่อ $B \cap C = \emptyset$
    \item ในการสำรวจนักศึกษา 120 คน พบว่าชอบเขียนโปรแกรม 75 คน ชอบระบบเครือข่าย 60 คน ถ้าทุกคนชอบอย่างน้อยหนึ่งอย่าง จงหาจำนวนคนที่ชอบทั้งสองอย่าง
    \item \textbf{(ประยุกต์คอมพิวเตอร์)} กำหนด $U = \{1, 2, \ldots, 8\}$, $A = \{1, 3, 5, 7\}$, $B = \{3, 4, 5, 6\}$
          จงเขียนเวกเตอร์บิตของ $A, B, A \cup B, A \cap B, A - B$ และตรวจสอบว่าการดำเนินการระดับบิต (OR, AND, AND NOT) ให้ผลลัพธ์ตรงกับทฤษฎีเซต
  \end{enumerate}
\end{frame}

\begin{frame}[plain]
  \vfill
  \centering
  {\Large\color{navy}\textbf{จบบทที่ 2 : ทฤษฎีเซต (Set Theory)}}\par
  \vspace{0.8em}
  {\large\color{Tgreen}\textbf{เตรียมพบกับ บทที่ 3 : ความสัมพันธ์และฟังก์ชัน}}\par
  \vspace{0.4em}
  {\color{gray} ซึ่งพัฒนาต่อยอดจากผลคูณคาร์ทีเซียน $A \times B$ ที่เราเรียนในวันนี้!}
  \vfill
\end{frame}
